\chapter[Myths, Worlds, and Lives]{Do Myths Explain the World or Arrange Our Lives?}
\section{The Moon inside a Mooncake}
Begin with an object: a mooncake.

Nothing rare. A month before Mid-Autumn, tins, snow-skin cakes, and flowing custard varieties occupy supermarket entrances in mountains. Yours is traditional Cantonese style, baked dark gold, moulded with a rabbit, tree, and woman in fluttering ribbons. Nobody needs to identify Chang'e, the Jade Rabbit, and lunar cassia tree.

Answer seriously: why eat it?

Because it tastes good is insufficient. Identical ingredients are available through the other eleven months, yet nobody thinks of mooncakes in March or July. On the fifteenth day of the eighth lunar month, the calendar needs this pastry more than your tongue. More precisely, an old story does: a woman swallowed an immortality elixir, flew to the moon, and remained there. For this unseen event, over a billion people look up on one night, share round cakes, and message distant family with festival wishes.

Notice the cake's work. It supplies no lunar facts, nothing of gravity, craters, tidal locking. Instead it selects one day from 365, stamps it special, and specifies tonight's companions, food, gaze, and memories.

Is its woman an obsolete astronomical report or a still-operative life schedule? \textbf{Do myths explain worlds or arrange lives?}

\section{A Rooftop Question}
Come into a scene.

Eight in the evening on the eighth lunar month's fifteenth day, atop an apartment block in a southern city. Daytime heat lingers; wind between buildings carries cooking oil. The moon has cleared the opposite roof, almost unreal in size.

A folding table holds Grandma's sliced mooncakes, peeled pomelo, and small bowl of water whose use she cannot explain, only remembering her mother placed one too. Dad positions a phone on the railing for moon photographs, Mum competes for digital red envelopes in a group chat, and thirteen-year-old An'an bites cake before speaking.

``Grandma, is Chang'e really on the moon?''

Distributing pomelo without looking up, she says yes: the dark patch is the cassia tree.

``Chang'e 6 brought far-side soil back last year,'' he says. ``The whole moon has been photographed countless times. No palace, no rabbit, only rock and dust. A rock.''

Dad puts down his phone, interested. ``Ancient people did not know what the moon was, so invented an explanation. Myth was their science. Now we understand and no longer need it.''

Grandma presses pomelo into An'an's hand. ``Stories are stories, but festivals still need celebrating. Without Chang'e, who would join an old woman on the roof tonight?''

Mum looks up from red-envelope competition. ``Exactly. When your dad was little, I had to make him memorise `Moonlight before my bed.'''

An'an bites the fruit, silent but unconvinced. Either something inhabits the moon or not. If proved false, how can the festival still be necessary? Is that not contradiction?

Remember this rooftop. Grandma and Dad use myth for entirely different things. We shall separate them, then examine their entanglement.

\section{Is Myth Obsolete Science?}
Illuminate the conflict before deciding a winner.

Dad's camp declares \textbf{myth primitive science}. Without electrical knowledge, thunder becomes a god; without bacterial knowledge, fever and delirium become possession; without knowing the moon as orbiting rock, its phases invite Chang'e. Lacking instruments and equations, people explain through stories. Scientific light retires them as electric lamps retire candles.

Clean, plausible, pleasantly progressive: we understand, they did not.

Grandma declines that contest. She never argues for a lunar palace and might agree about rock and dust. Her concern is why cakes, pomelo, and family gather here. Remove Chang'e and lunar data remain unchanged, but next year's folding table may never appear.

See the mismatch? Father and son, or rather Dad and Grandma, are not disputing one thing:

\begin{itemize}
\item Dad asks: \textbf{are the story's claims about the world true?} An explanatory question.
\item Grandma defends: \textbf{are the actions it brings about still worth doing?} A question about life.
\end{itemize}
Does a refuted myth die? Or was its work always arranging life rather than answering why? If both, how do the jobs relate? Why did every civilisation, rice-growing, pastoral, maritime, produce mountains of creation, flood, hero, and afterlife stories, so similar yet different?

Choose a rooftop position first. Is An'an's unease justified? Can you accept a demonstrably false story confidently scheduling the year?

\section{Two Truths on the Moon}
Strengthen both positions, then see how others handle the difficulty.

\textbf{First: Myth is ancient science. Give it its full case.}

Too often dismissed as disrespecting tradition, it deserves a hearing.

First, explanatory work really exists. Zeus's thunderbolts, Thor's sky-crossing chariot, Chinese thunder and lightning deities answer what something is. Many myths explicitly explain snakes, mortality, sunrise. Australian Aboriginal Dreaming stories correspond closely to terrain, water, animal behaviour, functioning as maps and survival manuals. This is not later overinterpretation.

Second, early science at least treats ancient people as serious thinkers. Nineteenth-century anthropologists Edward Tylor and James Frazer took such approaches. Frazer's \textit{The Golden Bough} arranged thought as magic, religion, science, each a rough predecessor. This progress ladder has major problems, shortly discussed, but retains respect for serious attempts with available materials to answer genuine questions, not idle tales for children. In that sense, myth precedes science rather than opposes it.

Third, science visibly displaced explanatory jobs. Nobody now relies on Chinese thunder deities to explain lightning, nor forecasts on sacrifice. Explanations changing generations are historical facts.

\textbf{Second: Myths were never only answers; they were schedules.}

Grandma's side is equally strong.

First, ancient people need not have believed as literally as we suppose. Treating myths as their unquestioned textbooks is an easy mistake. Cicero's \textit{On Divination} jokes, broadly, that two diviners meeting must struggle not to laugh at one another, despite daily predicting affairs through birds and livers. Insiders already doubted and joked two millennia ago. Israel's prophets sharply mock foreign images. Confucius answers spirits and death by prioritising service to living people. Does that sound simply literal?

Second, belief has varieties. Hawaiian Pele lives in a crater and erupts lava in anger. Some contemporary geologists also inherit her story tradition, measuring magma by day and chanting praise at festivals. No divided personality, but divided labour between truths: predicting eruptions versus knowing who we are and our relation to the island. Calling the latter a poor imitation of the former confuses occupations.

Third, myths organise calendars, bodies, relationships, identities, not only minds. Spring Festival, Dragon Boat Festival, Mid-Autumn carry stories. Jewish Passover dinners annually reenact Exodus, the youngest child asking four questions. Mexico's Day of the Dead combines altars, sugar skulls, returned ancestors, with roots in precolonial Mesoamerican death traditions. Nobody at Passover demands proof of the parted sea before eating. Stories bind us collectively, renewing the bond yearly.

Dad prosecutes the explanatory job; Grandma guards the life-arranging job. Historically, one story may have worked both shifts.

\section[Thinking Bridge: Take Myths Apart]{Thinking Bridge: Dismantle Myths into Components}
Can we analyse broad impressions, universal floods and heroes, through more than feeling? Use the portable tool \textbf{motif}.

A motif is a detachable, reusable story component: a warned person, boat or gourd, surviving pair, scouting bird, siblings repopulating humanity, hero returning from the underworld. Like building blocks in different castles, one appears thousands of kilometres apart.

This is scholarly method, not literary pastime. Early twentieth-century Finnish scholars developed the historical-geographical method: record versions with times and places, compare added and removed components, track journeys like viral mutations. American Stith Thompson extended it into the multivolume \textit{Motif-Index of Folk-Literature}, numbering thousands of elements, from descending creators to fire-winning animal helpers, enabling international cross-reference.

Practise on five distant flood traditions, comparing components:

\noindent
\begin{tabularx}{\linewidth}{llllX}
\toprule
Source & \shortstack[l]{Who is\\ warned} & Survival & \shortstack[l]{Nature\\ of flood} & Afterwards \\
\midrule
\shortstack[l]{Mesopotamia,\\ Gilgamesh} & \shortstack[l]{Utna-\\ pishtim,\\ secretly\\ warned\\ by a god} & \shortstack[l]{Large boat\\ with\\ animals} & \shortstack[l]{Gods\\ destroy\\ the world} & Birds scout; sacrifice ashore; immortality \\
\shortstack[l]{Hebrew\\ tradition,\\ Genesis} & \shortstack[l]{Noah,\\ instructed\\ by God} & \shortstack[l]{Ark with\\ paired\\ animals} & \shortstack[l]{Punishment\\ for human\\ sin} & Dove scouts; covenant and rainbow \\
\shortstack[l]{Ancient\\ Greece} & \shortstack[l]{Deucalion\\ and wife,\\ warned by\\ Prometheus} & \shortstack[l]{Wooden\\ chest afloat\\ nine days} & \shortstack[l]{Zeus\\ destroys\\ the world} & Stones thrown behind become people \\
% EN-PAGINATION-BEGIN flood-comparison
\bottomrule
\end{tabularx}

\noindent
\begin{tabularx}{\linewidth}{llllX}
\toprule
Source & \shortstack[l]{Who is\\ warned} & Survival & \shortstack[l]{Nature\\ of flood} & Afterwards \\
\midrule
% EN-PAGINATION-END flood-comparison
\shortstack[l]{Ancient\\ India,\\ Shatapatha\\ Brahmana} & \shortstack[l]{Manu,\\ warned by\\ a divine\\ fish} & \shortstack[l]{Boat tied\\ to fish's\\ horn, towed\\ uphill} & \shortstack[l]{All-\\ engulfing\\ waters} & Manu alone produces humanity \\
\shortstack[l]{Southern\\ Chinese\\ Nuwa and\\ Fuxi\\ traditions} & \shortstack[l]{Siblings,\\ warned by\\ Thunder\\ God or fore-\\ knowledge} & \shortstack[l]{Shelter in a\\ great gourd} & \shortstack[l]{World-\\ engulfing\\ waters} & Siblings marry and repopulate humanity \\
\bottomrule
\end{tabularx}

\smallskip
Know the table's abilities and limits. Cell-by-cell comparison replaces vague resemblance with visible similarities and differences. It cannot declare all versions descendants of one story. Similarity starts inquiry, not ends it; we shall return to this in comparing explanations.

Practise with Nezha: lotus incarnation, self-killing to repay parents, sea turmoil and dragon killing. Do these have relatives elsewhere? Death, transformation, rebirth recurs in Nezha, Egypt's murdered and dismembered Osiris resurrected as underworld ruler, Greece's Persephone whose annual crossings produce seasons. The motif chart maps mythology: neighbours, distant kin, and merely similar strangers become distinguishable.

\section{Reading a Two-Thousand-Year-Old Text}
Set paraphrase aside briefly. The \textit{Classic of Mountains and Seas}, ``Regions beyond the Seas: North,'' gives all of Kuafu's story in forty-odd characters:

\begin{quote}
Kuafu raced the sun and entered its light. Thirsting, he drank from the Yellow and Wei Rivers. They were insufficient; he went north to the Great Marsh. Before arriving, he died of thirst on the road. His discarded staff became Deng Forest.
\end{quote}
In modern terms: he chased into sunlight, drank the two rivers dry, sought more water northward, died before reaching it, and left a staff transformed into peach trees.

Study this specimen for a minute. Does it explain why? Barely: a forest's origin, casually supplied in the ending. What stays is someone daring to race and catch the sun, then dying, yet leaving shade for future people. Explanation or stance, astronomy or a short poem about pursuing what cannot be caught?

The \textit{Huainanzi}'s ``Surveying the Obscure'' tells Nuwa's sky repair:

\begin{quote}
In ancient times the four supports failed, the nine regions split, heaven no longer covered all, earth no longer bore all. \,\ldots\, Nuwa smelted five-coloured stones to mend the sky and cut a giant turtle's feet to restore the four supports.
\end{quote}
Broadly, fallen sky-pillars and cracked earth disrupted coverage and support; she repaired heaven with coloured stones and replaced pillars with turtle legs.

An honest qualification: Pangu separating heaven and earth may seem China's oldest story, yet its earliest written record is late, in the Three Kingdoms \textit{Sanwu Liji}, younger than many \textit{Mountains and Seas} passages. Creation need not open a tradition. A civilisation can tell floods, heroes, repairs for centuries before adding a world's beginning. A story's age differs from the age of its events.

Together, these texts show Chinese heroes often \textbf{repair} rather than defeat nature: mend leaking sky, manage floodwater, shoot nine excess suns while leaving one, as Hou Yi does. Different component lists are evidence for the next section.

Finally, modern evidence. Mesopotamian flood stories slept on buried tablets for over two millennia. In 1872, British Museum scholar George Smith deciphered fragments from Nineveh, Tablet Eleven of \textit{Gilgamesh}: Utnapishtim recounts divine flood, boat rescue, scouting birds, closely paralleling Genesis. Smith supposedly jumped with excitement. Europe was astonished that a tablet predating biblical composition told the same flood. Debate still burns over borrowing or shared real disaster.

Three pieces lie before us: poetic sun-chasing, repair-like sky-mending, sensational clay. Compare explanations.

\section{The World's Floods: Three Explanations}
Flood stories across Mesopotamia, India, Greece, southern China, and Indigenous American oral traditions are factual. Texts also record participants' understandings. Debate concerns \textbf{why so many exist}. Each account has reasons and limits.

\textbf{First: Stories travel (diffusion).}

Some resemblances exceed easy coincidence. Utnapishtim and Noah share not merely water but warning, boat, animals, mountain grounding, birds, sacrifice. Independently inventing that sequence resembles strangers dreaming the same television series. Most scholars therefore connect Hebrew flood narrative with earlier Mesopotamian tradition. No embarrassment: merchants, captives, migrants formed an exchange network, carrying stories.

Limit: neighbouring likeness is easier than distant likeness. Pre-navigation-era Americas were ocean-separated; Miao and Yao gourd stories differ greatly in components. Universal diffusion imagines one storytelling ancestor and everyone else mere messenger, without evidence and belittling creativity.

\textbf{Second: Situations create stories (shared predicaments).}

Early farming civilisations clustered beside great rivers: Mesopotamian, Nile, Indus, Yellow. Water sustains crops or destroys villages. Total flood requires memory, not imagination. Elders recount disastrous years; generations compress them into one world-drowning event. Likewise mortality and heroes seeking escape arise everywhere because death needs no teacher.

Limit: this explains floods, not identical details such as scouting birds, sometimes the same species. Shared situations are also so broad they can absorb almost any resemblance, risking explanation of everything and therefore nothing.

\textbf{Third: A great flood really happened (catastrophe memory).}

The boldest, most tempting account: after the Ice Age, seas rose over a hundred metres, drowning vast coasts in genuine global inundation across ten millennia. A specific geological hypothesis proposes a freshwater Black Sea flooded by Mediterranean water over seven thousand years ago, scattering survivors and memories. Chinese researchers also suggest an upper Yellow River flood around four thousand years ago, near Yu's traditional date. If established, myths become humanity's oldest disaster archive.

Limit: none is universally settled. Black Sea speed and scale are disputed; linking Yellow River evidence to Yu seems fragile to many. Even a real event passes through generations of retelling into a reshaped story. Disaster memory offers possible chronology, not explanation of narrative form.

\textbf{An easily misjudged counterexample: Do not let the table deceive.}

Five aligned stories feel satisfyingly related. Cold water: Chinese flood-control stories and Noah tell \textbf{nearly opposite things}. The ark centres on irresistible divine judgement and survival through warning and shelter. Yu treats water as engineering, not judgement. His father Gun's blocking fails; Yu drains, works thirteen years, passes home thrice without entering, and overcomes flood. One teaches survival through obedience, another transforming the world. Compare only water and survivors, and erase the difference.

Neat cells can uniform unlike souls. Motif charts discover questions; they do not announce answers.

\section[Further Challenge: Three Myth Lenses]{Further Challenge: Three Microscopes for Myth}
Now three twentieth-century theoretical microscopes asking what myth does, with incompatible answers worth comparing. Each was once treated as sole truth by some, then exposed for blind spots. Tools, not standard verdicts.

\textbf{First: Functionalism---myth as society's instrument.}

Stranded in the Trobriands during World War I, Bronislaw Malinowski stayed for years of fieldwork. His conclusion: myths are not idle fairy tales but \textbf{social charters}. An ancestral story validates family land; a myth prescribes fishing rites. Like a living constitution, narrative supplies antiquity's endorsement for power and custom. Its truth matters less than \textbf{whether it works}.

Durkheim reached similar territory through Australian Aboriginal rituals. People worship disguised society. Shared songs and movements generate a force exceeding individuals, named divine but actually collective cohesion. Ritual periodically recharges society.

Functionalism shifts priority from truth to action. Grandma's rooftop question is plain functionalism. Its limit: if myths sustain the status quo, whence rebellion, tragedy, irony? Which norm does Kuafu maintain? The instrument explains usefulness, not moving power.

\textbf{Second: Structuralism---myth as mental exercise.}

Claude Levi-Strauss saw neither primitive science nor merely social glue, but \textbf{the mind's handling of contradictions}: life/death, nature/culture, raw/cooked, human/divine. Unsolvable oppositions are rehearsed narratively. His Oedipus analysis disregards moving plot for component tables, symmetries, reversals, solving myth like equations.

Similarities might arise neither through travel nor flood, but because \textbf{humans share the same kind of brain} confronting irresolvable problems. Heroes returning from below rehearse whether death can be traversed, an unavoidable computation.

Its blind spot is conjuring. Different analysts derive different deep structures from one text with little means of deciding. The stronger an unfalsifiable explanation, the more caution it deserves, a lesson applicable to many theories here.

\textbf{Third: Psychological explanation---myth as inward projection.}

Carl Jung proposed a collective unconscious beyond individual memory: innate archetypes, hero, mother, sage, trickster, recurring through myths and dreams like actors touring theatres. Influenced by him, Joseph Campbell's \textit{The Hero with a Thousand Faces} describes a shared journey: departure, trials, treasure or wisdom, transformative return. The model entered Hollywood manuals; \textit{Star Wars} openly acknowledges its influence.

Honesty requires: \textbf{the collective unconscious is a hypothesis, not a verified discovery}. Mainstream psychology does not accept its literal form; no evidence establishes a preloaded myth-component library. The hero's journey works as screenplay formula, not proof of innate mental structure. Perhaps it is broad enough to fit anything. Distinguish usefulness from truth and this chapter has earned its place.

Three microscopes show social maintenance, cognitive exercise, inward projection. Not wholly exclusive: Chang'e schedules reunion for functionalists, stages immortality versus separation for structuralists, imagines losing a beloved for psychologists. The right lens depends on the question. How questions determine visibility may be our most valuable lesson.

\section{What Do Myths Wear Today?}
Myth is not dead antiquity. Look around: it changes clothes frequently.

\textbf{National dress.} Almost every state retells founding days, people, words in films and textbooks. Shared stories sustain communities, as Passover does; not inherently bad. Functionalist questions ask who stars, who disappears, which current arrangements gain endorsement. Asking is not disloyalty but assuming ownership of common stories.

\textbf{Consumer dress.} Nike names a victory goddess; Starbucks uses a sailor-luring siren, apt for supposedly irresistible coffee. Cars, perfumes, teams borrow mythic names. Advertising borrows structures: product as sacred object, purchase as initiation, improved self as reward, a shopping-centre hero's journey. Consumerism may be today's most successful ritual system, with festivals, mall temples, new-model holy objects, yearly reincarnation. Seeing it need not breed cynicism. Next time a counter quickens your pulse, recognise an ancient force.

\textbf{Technological dress.} China's Chang'e lunar programme, Jade Rabbit rover, Kuafu-1 solar satellite, Zhurong Mars rover send ancient stories into space to verify where myths fail. Rockets answer An'an. Yet stories persist. Singularity narratives, AI ending old rules and bringing elevation or extinction, resemble apocalypse and saviour: mighty arrival, old world's end, believers saved. Mars settlement recasts departure, trial, wilderness homeland in a spacesuit. Structural likeness does not prove scams; technologies are real. \textbf{Even future narratives wear ancient structures}, which quietly shape our wagers.

\textbf{Your clothes too.} Lucky-carp forwarding, unchanged exam avatars, lucky numbers are miniature magic, sharing the mechanisms Frazer recorded: uncertainty invites a rite of doing something. Ancient people knew it as ritual; you call it seeking luck.

Return to the roof. Dad is right that myths' explanatory posts retire, even to the moon's far side. Grandma is right that their organising posts expand into flags, logos, launchpads, shopping carts. Myth did not lose to us; it works without its old job title.

\section[Your Turn: Can We Live without Stories?]{Your Turn: Is a World without Stories Habitable?}
Return to cake and rooftop.

Imagine an efficiency planner banning unverified public stories: retain astronomical festival dates but erase legends, forbid mythic branding, allow only archival founding claims, punish lucky-carp forwarding as misinformation. Science explains; timetables organise.

Would this world be clearer or harder to inhabit?

Weigh your components before answering.

Primitive-science arguments recognise sincere explanatory work now largely replaced. Remember Frazer's ladder: depicting predecessors as uncivilised versions of us is arrogant and often false; their beliefs may not have been literal.

Functionalism shows that a false story can produce a real us and a real annual table. Factual and relational truth differ.

Motif charts and warnings reveal travel, shared situations, shared minds as possible resemblance sources, while tidy tables mistake distinct souls for copies. Yu and Noah prove the danger.

Three microscopes accompany the caution that theories are useful tools, not final judgements, and collective unconscious remains hypothetical.

Is a society without shared unverified stories---Mid-Autumn without Chang'e, nations without myths, technology without narratives---more honest or more impoverished? If stories are necessary for gathering, how should we select, revise, and tell them fairly to believers, omitted people, and future discoverers of truth?

No standard answer. But next eighth-month full moon, accepting Grandma's pomelo, your mooncake may taste different.
